\chapterpage{pipeline}{08 / THE CONCRETE PIPELINE}{Several verifiable deliveries.}{The public user path is the acceptance subject. A candidate on a PR branch may be well developed and tested without already affecting the delivered product.}
\par
{\tighttable
\begin{tabularx}{\textwidth}{@{}p{31mm}p{18mm}Xp{49mm}@{}}\toprule
Deliverable & Path & State read by the source report & Still-required effect\\\midrule
Mesh recovery and seed discovery & \#437 & Head \texttt{3c9c5004...}; 400-byte validation and native tests documented. & Main adoption, permitted Authority capability and productive recovery readback.\\
Personal Firefox / terminal & \#438 & Windows ARM64 client and server AMD64 controls documented; Win11 AMD64 remains HOLD. & Actual supported x64 client, public product path and acceptance.\\
Interruption-resilient work & \#445 & Contract preserved; actual client interruption test not yet executed. & Productive handoff, interruption/restart, continuation without duplicate effect.\\
Linux scaling & \#443 & Bounded 1/2/4-process speedup documented; cause of earlier I/O outliers open. & Controlled isolation for the general scaling claim.\\
Authority/Mirror A/B & \#447 & Diagnostic run: credential-delivery HOLD; no live speedup. & Access, byte/mode-equivalent method and full distributed trial.\\
Dependency replacement & \#448 & Standing Owner Directive and learning contract; green tests documented. & Concrete replacement that removes the original dependency.\\
Lifecycle governance & \#449 & Direct Main writer removed in candidate; separate promotion planned. & Effective Main protection, independent review and first actual lifecycle acceptance.\\
Scientific publication & \#447 & Existing small PDF package; new v3 contract proposal, not activated for production. & Contract decision, exact new fileset, return, single-use decision, upload/readback.\\\bottomrule
\end{tabularx}}
\vspace{2mm}{\sffamily\footnotesize Table 2. Source-supported work fronts, not a claim of complete inventory of all nodes or all Authority PRs. PR descriptions also contain historical partial states; the metadata head actually read is authoritative for that read. \src{R2,R4--R9}}
\begin{cols}
\subsection*{The directly visible benefit}
A first usable QIK-VRT workplace need not satisfy the whole hardware objective simultaneously. It must work within its clearly bounded scope: personal access, a persistently stored task, reliable resumption, a permitted actor and visible evidence of the completed task. This partial acceptance must not be called complete closure of the Cloud/Universal-Transputer and FPGA programme.

The personal browser provides access; the repository carries state-bound context; the terminal mediates tasks and readbacks. OAuth/SSO and biometrics belong to the desired product profile. The source evidence does not contain complete end-to-end acceptance of this authentication chain. A login concept must therefore not be presented as a completed biometric product binding.

\subsection*{A concrete startup dependency}
On Main \texttt{8153c4e6...}, read by the source report, \code{RUNTIME_DEPENDENCY_MANIFEST.json} still requires a Windows x64 Python runtime, an active download-consent path and \code{bundled_windows_python_runtime=false}. The document is classified as mutable, historically generated runtime state. It establishes its manifest content, not inspection of every possible delivery archive. \src{R12}

For a reliable offline or self-provisioned start, the existing runtime/cache path must bind exact payload, hash, license, provenance, restore recipe and self-test. The path is accepted only when the declared environment starts without an unbound surprise dependency. General permission to assimilate does not replace this individual check.

\subsection*{Governance is part of delivery}
In the source report, PR \#449 still documents missing effective Main protection and no independent Code-Owner review. The separately read review-execution status on that PR-447 head also failed although CI and other workflows could succeed. There is no contradiction: an observer can fulfill its observation task while finding missing approval. \src{R9,R13}

\subsection*{Hardware remains a separate milestone}
VHDL, synthesis, bitstream, programming and physical measurement are different acceptance steps. The supplied history records open board, programmer and physical readbacks. This edition observes no new board run. Neither a hardware date calculated from software CI nor a successful physical signal witness is claimed. \src{L4}
\end{cols}

\chapterpage{execution}{09 / FROM TASK TO DELIVERY}{The sequence follows blockers.}{The delta principle should reduce needless repetition. It removes neither global tests nor external dependencies. The next productive action must address the actual bottleneck.}
\begin{cols}
\subsection*{P0: Close shared bottlenecks}
First establish the authorized Authority reading path and effective repository governance. They remain two separate subjects. A token permitted to read source bytes does not automatically grant administrative rights. Protected Main does not supply missing Authority credentials either.

The concrete access check is: an already authorized credential is actually delivered through the existing bounded path to the correct runner; current Authority head and tree are then read and bound to source, run and role. The diagnostic result in the source shows only that delivery is not established in the observed process. It proves neither general absence of credentials nor the cause of a public 404 response. \src{R3}

Governance requires observed protection, an appropriate independent review decision and current gates. Only then is separate, expected-head-bound adoption of a candidate admissible. General Owner agreement does not replace the technical check of that particular transition. \src{R9,R13}

\subsection*{P1: Close one complete user path}
Next execute a finite pilot task through personal access: record the task, admit the action, use a program or service, save the result, interrupt, restart and continue the same task correctly. Every required tool and runtime must be bound in the existing cache/provision contract. A client must be able to reconstruct the result without another improvised dialogue step. \src{R1,R5,R12}

For Windows x64 this test takes place on the actually required Windows 11 client. A successful ARM64 path and server x64 run are useful evidence but not substitute platforms. A separately bounded Linux pilot may be possible earlier; its scope must remain explicit relative to missing Windows or hardware acceptance. \src{R6}

\subsection*{P2: Optimize after full correctness}
Correct the measurement method according to the audit. Then compare identical tasks without and with delta reuse. Mandatory gates stay equal; result quality and error rate must not deteriorate. Adopt a speed change as a benefit only when supported within the agreed measurement window and regression budget.

An optimization can be slower yet detect errors better. Quality and time objectives must therefore be reported separately. Conversely, a fast run without full readback is not success.

\subsection*{P3: Make findings broadly usable}
Only checked changes are integrated into existing runtime, handoff and publication paths. A successor receives its own subject binding. Older PRs are assessed against their intended effect; a subject already completed elsewhere can be classified as superseded, but not solely because of age.

The task ``on every node'' requires a current finite list of authorized targets. For each target, record delivery identity, local adoption and readback. A rule for future nodes does not show that all current nodes have been reached. Neither a global scan nor unauthorized self-propagation follows.

\subsection*{Check what defines completion}
A global mandatory check can reevaluate the entire relevant state despite a small delta. It is not algorithmic waste when required for the current claim. Repeated navigation, needless reconstruction of stable relationships and duplicates of the same work unit are avoidable instead.

\begin{boundary}{What the source edition actually applied}
Existing sources and layout methods were reused; old and current heads were separated; calculations were checked; two methodological errors were locally demonstrated; access, governance and measurement repair were added to the forecast. No production deployment, Main merge or unmeasured performance gain follows. This English successor adds translation and new artifact validation, not new historical trials.
\end{boundary}
\end{cols}

\chapterpage{forecast}{10 / THE TIME FORECAST}{A first delivery may be near. Overall completion has no date.}{A date is meaningful only when the product, acceptance scope and prerequisites actually available by then are clear.}
\begin{cols}
\subsection*{Changes from the earlier forecast in the source report}
The earlier CI-admission blocker at \texttt{e7e8737f...} no longer describes the aggregate state reported by the source. PR \#447 is at \texttt{a31ebb7e...}; its CI and A/B diagnostic ran. PR \#443 carries bounded positive scaling evidence. PR \#437 clarified seed discovery, and \#449 has a governance successor with fresh tests. \src{R2--R4,R7,R9}

Publication is also more than ``one approval away'': the v2 contract has a documented hash/authorization transition boundary; the v3 proposal has not been activated for production. The live measurement path needs methodological work too. These points extend or structure the critical path even if individual tests are already green. \src{R10,R11,R14}

\subsection*{Why 2.5 times faster is not 2.5 times earlier}
Let $p$ be the proportion of total time affected by local speedup $s$. Without additional overhead, arithmetic gives
\[
 S_{\mathrm{gesamt}}=\frac{1}{(1-p)+p/s}.
\]
The subscript denotes overall speedup. Additional normalized effort $h$ adds to the denominator. For $s=2.5$ and $p=0.1$, only 1.064 times overall speedup or six percent time savings would be possible. For $p=0.5$, it would be 1.429 times or 30\%. These are calculation scenarios; the actual $p$ for the delivery pipeline was not measured. \src{C1}

Waiting for devices, access, review or external acceptance is not automatically shortened. Parallel work is also limited by runners, independent reviewers, single writers and the productive action path. Freely assuming unlimited capacity would make the forecast unusable.

\subsection*{A checkable scheduling model}
For work unit $i$, let $r_i$ be the earliest availability of prerequisites, $d_i$ estimated active duration and $P_i$ its predecessors. A planning model is:
\[
 EF_i=\max\!\bigl(r_i,\max_{j\in P_i}EF_j\bigr)+d_i.
\]
If no date exists for necessary access, $r_i$ is unknown. There is then no justified fixed completion date. This names a dependency rather than making a blanket claim that engineering is inherently unpredictable.

\subsection*{Assumptions of the following window}
The calendar on the next page is an \textbf{editorial engineering estimate}, not a forecast fitted to measurements. It assumes an available executor, reachable independent reviewer, required runners/devices, no newly discovered severe defects and at most one normal correction cycle per package.

$T_0$ is not the time of agreement in chat. It is the actually observed time at which the respective prerequisites become available. Neither daily throughput nor success probability is invented. ``Days'' means elapsed processing days with active work under this capacity assumption; inactivity shifts the calculation.

\begin{boundary}{The concrete answer within the historical scenario}
A narrowly bounded software pilot appears plannable over a few days to about a week under these assumptions. The evidence provides no responsible completion date for the entire system including SSO/biometrics, all nodes and physical FPGA acceptance.
\end{boundary}
\end{cols}
\clearpage\markboth{10 / Delivery windows}{}
{\sffamily\small\color{Accent}10 / CONDITIONAL CALENDAR PLANNING}\par\vspace{4mm}
{\fontsize{24}{29}\selectfont What can be accepted, and when.}\par\vspace{4mm}
\textbf{Calculation example:} Required prerequisites are actually available on \textbf{4 October 2026}. This premise was \textbf{not satisfied} in the source report's readback. These windows show consequences of that preserved scenario, not promised dates or new present-state observations.
\vspace{4mm}
\par
{\tighttable
\begin{tabularx}{\textwidth}{@{}L{35mm}L{26mm}L{29mm}>{\raggedright\arraybackslash}X@{}}\toprule
Milestone & Active time remaining & Example calendar & Acceptance / dependency\\\midrule
Bounded Authority read path & 0.5--1 day after actual delivery & 4--5 October & Credential reaches correct runner; current HEAD/TREE readback. No expansion of rights.\\
Corrected A/B trial & 1--3 days after access and test environment & 5--7 October & Equal byte/mode definitions, full timing boundary, controlled AB/BA pairs and accepted readback.\\
Governance/recovery candidate effective & 1--3 days after protection and review & 5--7 October & Separate promotion, new Main binding, first actual recovery/lifecycle effect.\\
Bound Windows runtime & 1--2 days after available payload and adapter path & 5--6 October & Version/hash/license/provenance, restore and startup test. No unplanned download.\\
Windows 11 AMD64 witness & 1--3 days after suitable client & 5--7 October & Native architecture, persistent state, interruption/restart, replay refusal and public readback.\\
Actual task continuity & 2--5 days after productive handoff & 6--9 October & Intervening question and client interruption without Owner restart; no duplicate external effect.\\
Bounded personal pilot & 3--7 days after all pilot prerequisites & 7--11 October & One continuous publicly reachable user path, documented platform scope and fresh acceptance.\\
Zenodo complete edition & 1--3 days after contract and fileset closure & Earliest 5--7 October in the full-readiness scenario & Activated valid publisher contract, new return, exact single-use decision, identical uploaded bytes and redownload.\\
IETF companion / arXiv & Preparation: days; acceptance open & No acceptance date & Validation and complete sources; actual submission separate.\\
SSO/biometrics, all nodes, complete FPGA objective & No supported remaining duration & Open & Missing implementation/carrier/integration and end-to-end evidence for each.\\\bottomrule
\end{tabularx}}
\vspace{5mm}
\begin{boundary}{Do not add incompatible time gains}
Local factors 2.5, 1.624 and 2.141 arise from different subjects. They are neither multiplied together nor transferred to the entire roadmap. The pilot calendar applies only to a predetermined scope. Unavailable access still has no automatic provisioning date.
\end{boundary}
\vspace{3mm}
The next supported forecast arises once actual approval, runner, handoff and measurement times are recorded. Observed durations then replace particular assumptions; remaining variation must be preserved. A negative trial may lengthen the estimate, a reproducibly smaller delta may shorten it.

\chapterpage{publication}{11 / PUBLICATION}{A publishable document is not yet a published result.}{Repository, Zenodo, IETF, arXiv and Wikipedia have different purposes and acceptance boundaries. They are not collapsed into one ``published'' status.}
\begin{cols}
\subsection*{The new scientific edition described by the source}
The complete edition is a new candidate with new bytes, extended methodological findings and a new forecast. It does not silently replace a previously authorized PDF. Text sources, calculation artifacts, claim register and change notice accompany the PDF. The layout template remains a design source, not substantive peer review. This English successor is separately identified again.

The edition carries no unfulfilled kernel classification. New formal statements specify model, assumptions and argument; no machine proof is claimed. Document QA checks bindings, calculations, coverage and presentation, not the truth of the entire external world.

\subsection*{The active Zenodo boundary in the source report}
The v2 policy requires exact candidate bytes, complete claim disposition, appropriate evidence, boundary tests, return delivery and then a candidate-specific single-use decision. After the effect, it requires public record verification, byte-identical redownload and repository persistence. A general ``approval granted'' does not replace the later exact-byte decision. \src{R10}

The v3 proposal prepared at \texttt{a31ebb7e...} separates immutable proof readiness from later upload authorization. It addresses the documented conflict in which the earlier bundle includes an authorization field whose change changes the hash to be authorized. The proposal leaves active v1/v2 contracts and the production publisher unchanged. It is expressly \emph{not activated}. The binding final appendix preserves the later, separately identified implementation disposition without rewriting this historical statement. \src{R14}

Thus no apparently final \code{AUTHORIZE_EXACT_UPLOAD} line is generated for the complete edition while the valid production contract, complete binding and actually executable path remain open. No token, authorization or upload is invented.

\subsection*{IETF: protocol semantics, not consciousness certification}
The public Datatracker readback in the source report lists individual Experimental draft \code{draft-lohmann-qikvrt-effect-ack-03}. Its existing core does not require another wire format solely because of a new philosophical interpretation. A separate epistemic-status profile is a useful working candidate but needs its own specification and validation. Internet-Draft, RFC and IETF consensus are distinct. \src{W1,L4}

\subsection*{arXiv and Wikipedia}
arXiv provides scientific submission and distribution. Registration, endorsement where required, complete source materials and moderation are separate prerequisites. Moderation is not peer review; no arXiv publication is claimed here. The English successor supplies the complete English text required by the documented language correction, without submission. \src{W6}

Wikipedia should not establish a new discovery through self-publication. The prepared path described by the source is a neutral draft with disclosed conflict of interest and assessment of independent coverage. That edition performed no comprehensive new notability search and claims neither established significance nor established insignificance. \src{L4}

\subsection*{Preservation is a technical obligation}
Versioned sources, multiple authorized storage locations, checksums and restore tests reduce loss risk. They do not logically guarantee that data can never be lost. A hash does not replace a copy; a copy does not replace a recovery test; a node-inheritance rule does not replace actual readback on every node.

\begin{boundary}{Document delivery, not platform completion}
The source edition supplies a research article, an accessible synthesis and checkable companion artifacts. This successor supplies their complete English manuscript translation with new bindings. Zenodo publication, Main integration, node replication and product release remain separate effects with their own evidence.
\end{boundary}
\end{cols}
